\documentclass[12pt,a4paper]{article}

\usepackage[margin=1in]{geometry}
\usepackage{amsmath, amssymb}
\usepackage{booktabs}
\usepackage{enumitem}
\usepackage{hyperref}
\usepackage{parskip}
\usepackage{titlesec}
\usepackage{array}
\usepackage{xcolor}
\usepackage{fancyhdr}
\usepackage{graphicx}

\hypersetup{
    colorlinks=true,
    linkcolor=blue!60!black,
    urlcolor=blue!60!black,
}

\pagestyle{fancy}
\fancyhf{}
\rhead{System Requirements and Design Specifications}
\lhead{QUBO-Based DFL Node Placement}
\cfoot{\thepage}

\titleformat{\section}{\large\bfseries}{}{0em}{}[\titlerule]
\titleformat{\subsection}{\normalsize\bfseries}{}{0em}{}

\title{
    \textbf{System Requirements and Design Specifications} \\[0.5em]
    \large QUBO-Based Sensor Node Placement for Device-Free Localization
}
\author{Khalil Brik}
\date{Spring 2026}

\begin{document}

\maketitle
\tableofcontents
\newpage

% ─────────────────────────────────────────────────────────────────────────────
\section{Introduction}

This document describes the system requirements and design specifications for
the QUBO-based sensor node placement framework developed as part of the senior
thesis project. The system addresses the problem of selecting which Access
Point (AP) and Monitoring Point (MP) pairs to activate in a Device-Free
Localization (DFL) environment, given fixed hardware budgets.

The core idea is to formulate the pair selection problem as a Quadratic
Unconstrained Binary Optimization (QUBO) problem and solve it using both
Simulated Annealing (SA) and Simulated Quantum Annealing (SQA) via the
OpenJij library. A greedy baseline is also implemented for comparison.
Performance is measured using mean localization error from a k-Nearest
Neighbor (k-NN) classifier evaluated on held-out test fingerprints.

% ─────────────────────────────────────────────────────────────────────────────
\section{System Overview}

\subsection{Purpose}

The purpose of the system is to identify the best subset of AP-MP sensor
links for indoor fingerprint-based localization. Activating all available
links is not always practical due to hardware, power, or cost constraints.
The system finds a budget-feasible subset that maximizes localization accuracy
while minimizing redundancy between selected links.

\subsection{Scope}

The system covers the full pipeline from raw fingerprint data to evaluated
node placement results. This includes data preprocessing, QUBO formulation,
optimization via annealing solvers, localization evaluation, and figure
generation for the thesis. It does not include real-time hardware deployment
or wireless communication protocols.

\subsection{Stakeholders}

\begin{itemize}[leftmargin=2em]
    \item \textbf{Thesis author:} develops, runs, and evaluates the system.
    \item \textbf{Thesis committee:} reviews the methodology and results.
    \item \textbf{Future researchers:} may extend the system to real hardware
          or larger-scale environments.
\end{itemize}

% ─────────────────────────────────────────────────────────────────────────────
\section{Functional Requirements}

\subsection{FR-1: Data Ingestion}

The system shall load the following precomputed input data from CSV files:

\begin{itemize}[leftmargin=2em]
    \item RSS fingerprint matrices for training and test sets, of shape
          $K_{\text{train}} \times L$ and $K_{\text{test}} \times L$
          respectively, where $K$ is the number of fingerprint locations and
          $L$ is the total number of candidate AP-MP pairs.
    \item 2D coordinate matrices for training and test fingerprint locations,
          each of shape $K \times 2$.
    \item A pairs table listing the AP index and MP index for each of the
          $L$ candidate links.
    \item An importance score vector of length $L$, where each entry
          represents the RSS variance of the corresponding link.
    \item A redundancy matrix of shape $L \times L$, where each entry
          represents the correlation between the RSS signals of two links.
\end{itemize}

\subsection{FR-2: QUBO Formulation}

The system shall construct a QUBO matrix $Q$ parameterized by a scalar
$\alpha \in [0, 1]$ as follows:

\begin{equation}
    Q_{pp} = -\alpha \cdot I_p, \quad p = 1, \ldots, L
\end{equation}
\begin{equation}
    Q_{pq} = 2(1 - \alpha) \cdot R_{pq}, \quad p < q
\end{equation}

where $I_p$ is the importance score of pair $p$ and $R_{pq}$ is the
redundancy between pairs $p$ and $q$. The diagonal terms reward selecting
important links, and the off-diagonal terms penalize selecting redundant ones.

\subsection{FR-3: Binary Search over Alpha}

The system shall find the optimal trade-off parameter $\alpha^{*}$ using
binary search over the interval $[0, 1]$ with tolerance $\epsilon = 0.01$.
At each iteration, the QUBO is solved and the resulting selection is checked
for feasibility. A selection is feasible if the number of distinct AP nodes
is at most $n = 5$ and the number of distinct MP nodes is at most $m = 5$.
The binary search moves the lower bound up when feasible and the upper bound
down when over budget.

\subsection{FR-4: Optimization via Annealing Solvers}

The system shall support two annealing-based solvers:

\begin{itemize}[leftmargin=2em]
    \item \textbf{Simulated Annealing (SA):} implemented via
          \texttt{openjij.SASampler} with 100 reads and 1000 sweeps per call.
    \item \textbf{Simulated Quantum Annealing (SQA):} implemented via
          \texttt{openjij.SQASampler} with 10 reads and 1000 sweeps per call,
          emulating the behavior of a D-Wave quantum annealer with a projected
          hardware anneal time of 20~$\mu$s per read.
\end{itemize}

Both solvers use a fixed random seed of 42 for reproducibility.

\subsection{FR-5: Greedy Baseline}

The system shall implement a deterministic greedy baseline that selects
AP-MP pairs in descending order of importance score while enforcing the same
$n = 5$ and $m = 5$ budget constraints. This baseline serves as a
performance lower bound for comparison against the QUBO-based methods.

\subsection{FR-6: Localization Evaluation}

The system shall evaluate any candidate selection using a k-Nearest Neighbor
(k-NN) localizer with $k = 3$. For each test fingerprint, the localizer
computes Euclidean distances to all training fingerprints using only the
selected subset of RSS measurements, identifies the three nearest neighbors,
and predicts the location as their mean coordinate. The evaluation metric is
the mean Euclidean error in meters over all test locations.

\subsection{FR-7: Time-to-Solution Metric}

The system shall compute the Time-to-Solution (TTS) for each solver at the
optimal $\alpha^{*}$. TTS is defined as the expected total annealing time
needed to find the optimal solution with probability $p_R = 0.99$:

\begin{equation}
    \text{TTS}(\tau, p_R) = \tau \cdot \frac{\ln(1 - p_R)}{\ln(1 - p_s)}
\end{equation}

where $\tau$ is the anneal time per single run and $p_s$ is the empirical
success probability estimated from the reads. For SA, $\tau$ is the measured
wall-clock time per read. For SQA, $\tau$ is the projected D-Wave hardware
anneal time of 20~$\mu$s.

\subsection{FR-8: Multi-Run Statistical Experiment}

The system shall repeat the full binary search pipeline 100 times with
different random seeds to produce mean, standard deviation, minimum, and
maximum statistics for each solver. The greedy baseline is deterministic and
is therefore run once and replicated across all rows for consistency.

\subsection{FR-9: Result Persistence}

The system shall save all outputs to the \texttt{data/output/} directory,
including the binary search log, selected pair lists for each solver, a
per-run multi-run results table, a statistical summary, and a greedy result
file.

\subsection{FR-10: Figure Generation}

The system shall generate the following publication-quality figures at 300~DPI
and save them to the \texttt{figures/} directory:

\begin{itemize}[leftmargin=2em]
    \item Mean localization accuracy vs. $\alpha$ for SA and SQA.
    \item Time-to-Solution vs. $\alpha$ for both solvers.
    \item Number of selected pairs vs. $\alpha$.
    \item Number of active APs and MPs vs. $\alpha$.
    \item Per-solver AP and MP counts at $\alpha^{*}$.
    \item Total computation time per solver.
    \item Empirical CDF of mean localization error across 100 runs.
    \item Coverage comparison illustrating clustered, parallel, and diverse
          link configurations.
    \item DFL deployment diagram.
    \item System architecture diagram.
\end{itemize}

% ─────────────────────────────────────────────────────────────────────────────
\section{Non-Functional Requirements}

\subsection{NFR-1: Reproducibility}

All stochastic components of the system shall use a fixed seed so that results
can be reproduced exactly from the saved data and scripts.

\subsection{NFR-2: Performance}

The full binary search for a single solver shall complete within a reasonable
time on a standard laptop (under 10 minutes for SA with the given parameter
settings). The multi-run experiment of 100 trials may run for longer and is
intended to be executed once as a batch job.

\subsection{NFR-3: Portability}

The system shall run on any machine with Python 3.8 or later and the
following packages installed: \texttt{numpy}, \texttt{pandas},
\texttt{matplotlib}, and \texttt{openjij}. No proprietary quantum hardware
is required since the quantum annealer is simulated by OpenJij.

\subsection{NFR-4: Modularity}

Each stage of the pipeline is implemented as a separate numbered script.
This allows individual stages to be re-run independently without reprocessing
the entire pipeline from scratch.

\subsection{NFR-5: Output Quality}

All figures shall be saved at 300~DPI in PNG format (and in PDF format where
applicable) using IEEE-style font settings (serif, 9~pt) to ensure they are
suitable for direct inclusion in the thesis document.

% ─────────────────────────────────────────────────────────────────────────────
\section{System Architecture}

\subsection{Component Overview}

The system is organized into four logical layers, each corresponding to a
subdirectory in the repository.

\begin{enumerate}[leftmargin=2em]
    \item \textbf{Data layer} (\texttt{data/input/}): holds the raw and
          precomputed CSV files that serve as inputs to the optimization and
          evaluation pipeline.
    \item \textbf{Computation layer} (\texttt{scripts/}): contains the Python
          scripts that perform QUBO construction, binary search, evaluation,
          and baseline comparison.
    \item \textbf{Output layer} (\texttt{data/output/}): stores all
          numerical results produced by the computation layer.
    \item \textbf{Visualization layer} (\texttt{figures/}): stores all
          figures generated from the output data.
\end{enumerate}

\subsection{Data Flow}

Figure~\ref{fig:dataflow} summarizes the data flow through the system.

\begin{figure}[h!]
\centering
\begin{tabular}{c}
\texttt{data/input/} \\
$\downarrow$ \\
\texttt{scripts/01\_qubo\_binary\_search\_solver.py} \\
$\downarrow$ \\
\texttt{data/output/} \\
$\downarrow$ \\
\texttt{scripts/02 through 07 (plot scripts)} \\
$\downarrow$ \\
\texttt{figures/}
\end{tabular}
\caption{High-level data flow through the system pipeline.}
\label{fig:dataflow}
\end{figure}

\subsection{Script Responsibilities}

\begin{table}[h!]
\centering
\renewcommand{\arraystretch}{1.3}
\begin{tabular}{@{}lp{9cm}@{}}
\toprule
\textbf{Script} & \textbf{Responsibility} \\
\midrule
\texttt{01\_qubo\_binary\_search\_solver.py}
    & Builds the QUBO matrix, runs binary search for SA and SQA,
      evaluates via k-NN, computes TTS, and saves all results. \\
\texttt{02\_plot\_result\_figures.py}
    & Reads the binary search log and summary to produce the main
      result figures (accuracy, TTS, pair counts, solver comparison). \\
\texttt{03\_greedy\_baseline.py}
    & Implements and evaluates the greedy baseline and prints a
      three-way comparison against SA and SQA. \\
\texttt{04\_multirun\_experiment.py}
    & Repeats the binary search 100 times with varying seeds and
      saves per-run and summary statistics. \\
\texttt{05\_plot\_coverage\_comparison.py}
    & Generates a three-panel figure illustrating clustered, parallel,
      and diverse link configurations with Fresnel zone overlays. \\
\texttt{06\_plot\_dfl\_deployment.py}
    & Generates the DFL deployment diagram showing AP and MP nodes
      in the environment. \\
\texttt{07\_plot\_system\_diagram.py}
    & Generates the system architecture diagram for the thesis. \\
\bottomrule
\end{tabular}
\caption{Script responsibilities in the computation layer.}
\end{table}

% ─────────────────────────────────────────────────────────────────────────────
\section{Key Design Decisions}

\subsection{QUBO Parameterization}

The single scalar $\alpha$ was chosen as the only tunable parameter of the
QUBO formulation. When $\alpha = 1$, the objective reduces to pure importance
maximization with no redundancy penalty. When $\alpha = 0$, it becomes pure
redundancy minimization. Binary search over $\alpha$ is an efficient way to
find the point where budget feasibility is just barely satisfied, which tends
to correspond to the most accurate selection.

\subsection{Binary Search Instead of Constraint Encoding}

Budget constraints are enforced through binary search rather than being
encoded as penalty terms in the QUBO. This avoids the need to tune penalty
weights, which is a common difficulty with constrained QUBO formulations.
At each binary search iteration, feasibility is simply checked by counting
the distinct AP and MP nodes in the best sample returned by the solver.

\subsection{OpenJij as the Solver Backend}

OpenJij was selected because it provides both SA and SQA samplers through a
unified interface that closely mirrors the D-Wave Ocean SDK. This makes it
straightforward to switch to real quantum hardware in future work without
changing the surrounding pipeline.

\subsection{k-NN as the Localization Evaluator}

A 3-NN classifier was used to evaluate localization quality because it is
parameter-light, interpretable, and consistent with prior work in
fingerprint-based indoor localization. The same evaluator is used for all
solvers so that comparisons are fair.

% ─────────────────────────────────────────────────────────────────────────────
\section{System Parameters Summary}

\begin{table}[h!]
\centering
\renewcommand{\arraystretch}{1.3}
\begin{tabular}{@{}lll@{}}
\toprule
\textbf{Parameter} & \textbf{Value} & \textbf{Description} \\
\midrule
$N$              & 10         & Total number of sensor nodes \\
$n$ (AP budget)  & 5          & Maximum distinct AP nodes allowed \\
$m$ (MP budget)  & 5          & Maximum distinct MP nodes allowed \\
$\epsilon$       & 0.01       & Binary search convergence tolerance \\
$k$              & 3          & Number of neighbors in k-NN localizer \\
$p_R$            & 0.99       & Target success probability for TTS \\
SA reads         & 100        & Number of reads per QUBO call (SA) \\
SA sweeps        & 1000       & Number of sweeps per read (SA) \\
SQA reads        & 10         & Number of reads per QUBO call (SQA) \\
SQA sweeps       & 1000       & Number of sweeps per read (SQA) \\
$\tau_{\text{QA}}$ & 20 $\mu$s & Projected D-Wave hardware anneal time \\
Seed             & 42         & Random seed for reproducibility \\
$N_{\text{runs}}$ & 100       & Number of trials in the multi-run experiment \\
\bottomrule
\end{tabular}
\caption{Summary of key system parameters.}
\end{table}

% ─────────────────────────────────────────────────────────────────────────────
\section{File and Directory Structure}

\begin{verbatim}
Final_Submission/
    README.md
    system_requirements_and_design.tex      <- this document
    notebooks/
        01_precompute_fingerprint_data.ipynb
        02_qubo_solver_exploration.ipynb
    scripts/
        01_qubo_binary_search_solver.py
        02_plot_result_figures.py
        03_greedy_baseline.py
        04_multirun_experiment.py
        05_plot_coverage_comparison.py
        06_plot_dfl_deployment.py
        07_plot_system_diagram.py
    latex/
        algorithm.tex
        qubo_formulation.tex
    data/
        input/
            fingerprint_rss_train.csv
            fingerprint_rss_test.csv
            fingerprint_coords_train.csv
            fingerprint_coords_test.csv
            ap_mp_pairs.csv
            importance_scores.csv
            redundancy_matrix.csv
        output/
            binary_search_log.csv
            selected_pairs_sa.csv
            selected_pairs_qa.csv
            greedy_result.csv
            multirun_results.csv
            multirun_summary.csv
            summary.csv
    figures/
        fig00_dfl_deployment.png
        fig01_accuracy.png / .pdf
        fig02_time_to_solution.png / .pdf
        fig03_pairs_vs_alpha.png / .pdf
        fig04_ap_mp_vs_alpha.png / .pdf
        fig05_ap_mp_per_solver.png
        fig06_total_computation_time.png
        fig07_empirical_cdf.png
        fig08_coverage_comparison.png
        fig09_system_architecture_diagram.png
        coverage_animation.html
\end{verbatim}

% ─────────────────────────────────────────────────────────────────────────────
\section{Dependencies}

\begin{table}[h!]
\centering
\renewcommand{\arraystretch}{1.3}
\begin{tabular}{@{}lll@{}}
\toprule
\textbf{Package} & \textbf{Version} & \textbf{Purpose} \\
\midrule
Python     & 3.8+   & Runtime environment \\
NumPy      & any    & Numerical arrays and linear algebra \\
pandas     & any    & CSV I/O and data manipulation \\
matplotlib & any    & Figure generation \\
openjij    & any    & SA and SQA QUBO samplers \\
\bottomrule
\end{tabular}
\caption{Software dependencies.}
\end{table}

Install all dependencies with:

\begin{verbatim}
pip install numpy pandas matplotlib openjij
\end{verbatim}

% ─────────────────────────────────────────────────────────────────────────────
\section{How to Run the System}

All scripts should be executed from the \texttt{Final\_Submission/} root
directory so that the relative paths \texttt{data/} and \texttt{figures/}
resolve correctly.

\begin{enumerate}[leftmargin=2em]
    \item Run \texttt{notebooks/01\_precompute\_fingerprint\_data.ipynb} in
          Jupyter to generate the input CSV files. Skip this step if the
          \texttt{data/input/} files are already present.
    \item Run \texttt{scripts/01\_qubo\_binary\_search\_solver.py} to perform
          the binary search for SA and SQA and save the results.
    \item Run \texttt{scripts/02\_plot\_result\_figures.py} to generate the
          main result figures.
    \item Run \texttt{scripts/03\_greedy\_baseline.py} to evaluate the greedy
          baseline and compare it against SA and SQA.
    \item Run \texttt{scripts/04\_multirun\_experiment.py} to run the 100-trial
          statistical experiment.
    \item Run scripts 05 through 07 in any order to generate the remaining
          figures.
\end{enumerate}

\end{document}
